% Source-derived chapter 01: Introduction
% Original source: canonical-representations-surface-groups
\chapter{Introduction}
\label{canonical-representations-surface-groups-chapter-01}
\source{canonical-representations-surface-groups}

\begin{remark}[Source section]
\label{canonical-representations-surface-groups-chapter-01-source}
\source{canonical-representations-surface-groups}
\source{canonical-representations-surface-groups}
This chapter reproduces the corresponding section of the retrieved
LaTeX source, including its definitions, statements, examples, and
proofs. It has not yet been translated into Lean.
\notready
\end{remark}

% The source section title was: Introduction
\label{section:introduction}
\source{canonical-representations-surface-groups}
\subsection{Overview}
\label{subsection:overview}
\source{canonical-representations-surface-groups}
Let $\Sigma_{g,n}$ be an orientable topological surface of genus $g$ with $n$ punctures, and let $x\in\Sigma_{g,n}$ be a basepoint.
We denote the mapping class group of $\Sigma_{g,n}$ by $\on{Mod}_{g,n}$. 
There is a natural outer action of $\on{Mod}_{g,n}$ on $\pi_1(\Sigma_{g,n},x)$, induced by the action of $\on{Homeo}^+(\Sigma_{g,n})$ on $\Sigma_{g,n}$. 
Hence, $\on{Mod}_{g,n}$ acts on the set of conjugacy classes of
representations of $\pi_1(\Sigma_{g,n},x)$ into $\on{GL}_r(\mathbb{C})$. Our goal is to study the finite orbits of this action
via Hodge-theoretic and arithmetic techniques. 

\begin{definition}
\source{canonical-representations-surface-groups}
\notready
	\label{definition:mcg-finite}
\source{canonical-representations-surface-groups}
	For $g,n,r \geq 0$,
	a representation $\rho: \pi_1(\Sigma_{g,n},x) \to \on{GL}_r(\mathbb C)$
	is {\em MCG-finite} if the conjugacy class of $\rho$ has finite orbit under the action of
	$\on{Mod}_{g,n}$.
\end{definition}
Note here that we are studying \emph{conjugacy classes} of representations. For
example, semisimple MCG-finite representations correspond to finite
$\on{Mod}_{g,n}$-orbits in the character variety of $\pi_1(\Sigma_{g,n})$, not
its representation variety. Indeed, $\on{Mod}_{g,n}$ does not naturally act on
the representation variety of $\pi_1(\Sigma_{g,n})$.

Any representation of $\pi_1(\Sigma_{g,n})$ constructed without making choices
(e.g.~without naming a specific curve in $\Sigma_{g,n}$),
or by making choices with only a finite amount of indeterminacy, is MCG-finite.
So we view MCG-finite representations 
as the \emph{canonical} representations of $\pi_1(\Sigma_{g,n})$. 

The study of MCG-finite
representations has a long history.
%albeit with different terminology, 
It has a close
connection to isomonodromy and the Painlev\'e VI equation,
originally introduced in 1902 by Painlev\'e 
\cite{painleve:on-the-differential-equations-of-sectond-order}
and Gambier
\cite{gambier:on-the-differential-equations-of-second-order}.
For each $g,n,r$, conjugacy classes of semisimple MCG-finite representations 
$$\pi_1(\Sigma_{g,n})\to \on{GL}_r(\mathbb{C})$$ correspond to algebraic solutions
to a certain \emph{isomonodromy} differential equation, see \cite[Theorem A]{CH:isomonodromic}. 
Algebraic solutions to the Painlev\'e VI equation in an appropriate choice of coordinates 
correspond to MCG-finite representations with trivial determinant in
the particular case $g=0, n=4, r=2$ (see
\cite[\S1-\S3]{mahoux:itroduction-to-the-theory-of-isomonodromic-deformations}
and \cite[\S2]{doran:isomonodromic}). 

There is substantial literature on the
search for algebraic solutions 
to the Painlev\'e VI equation,
\cite{andreev2002transformations, boalch2005klein, boalch2006fifty, boalchsome, boalch2007higher, cantat2009holomorphic, dubrovin1996geometry, dubrovin2000monodromy, hitchin1996poncelet, hitchin2003lecture, kitaev2006grothendieck, kitaev2005remarks},
culminating in the classification given in \cite{lisovyy2014algebraic}. Boalch \cite{boalch2007towards} pitches the problem of classifying algebraic solutions to Painlev\'e VI as a natural generalization of Schwarz's list of hypergeometric equations with finite monodromy \cite{schwarz1873ueber}.

We will later see that the classification of MCG-finite representations is
closely connected to several major open questions in low-dimensional topology.
In particular, it has connections to the Putman-Wieland conjecture
\cite[Conjecture 1.2]{putmanW:abelian-quotients} and Ivanov's conjecture that mapping class groups do not virtually surject onto $\mathbb{Z}$ \cite[\S7]{Ivanov:problems}.
These conjectures in turn have applications to algebraic geometry, as they would
determine the first rational homology of finite covers of $\mathscr M_{g,n}$.

Many interesting classes of
representations are MCG-finite.
Examples include 
the rigid local systems studied by Katz \cite{katz2016rigid}, 
representations constructed via TQFT techniques ~(e.g.~\cite{kuperberg2011denseness, koberda2016quotients, koberda2018irreducibility, koberda2018representation}),
and representations of algebro-geometric interest, such as those constructed via the Kodaira-Parshin trick (e.g.~\cite[Example 3.3.1]{lawrence2019representations}).
We will see later that MCG-finite representations of $\pi_1(\Sigma_{g,n})$ are closely related to representations of finite-index subgroups of $\on{Mod}_{g,n+1}$. 
%In \autoref{section:questions} we speculate on links between MCG-finite representations and representations of geometric origin.



\subsection{Results for surface groups}
\label{subsection:surface-groups}
\source{canonical-representations-surface-groups}
Given the difficulty of classifying MCG-finite representations in the $g = 0, n = 4, r = 2$ case (encompassing Painlev\'e VI), it may be surprising that
 we are able to obtain a complete and simple characterization of MCG-finite representations when 
$g \gg r$.
%We next describe a complete and simple characterization of MCG-finite representations when 
%$g$ is large.
%That such a simple classification exists may be particularly surprising, given the
%difficulty of classifying MCG-finite reps in the $g = 0, n = 4, r = 2$ case, of which a special case
%is classifying algebraic solutions to the Painlev\'e VI equation.

\begin{theorem}
\source{canonical-representations-surface-groups}
\notready\label{theorem:finite-image}
\source{canonical-representations-surface-groups}
	For $g, n, r\geq 0$, let
	$$\rho: \pi_1(\Sigma_{g,n})\to \on{GL}_r(\mathbb{C})$$ be a MCG-finite
representation. If $r<\sqrt{g+1}$, then $\rho$ has finite image. 
\end{theorem}
	This answers a question of Junho Peter Whang \cite[Question
	1.5.3]{lawrence2019representations}. We will prove
	\autoref{theorem:finite-image} in
\autoref{subsection:main-proof}.
	
	While the statement appears to be purely topological or even group-theoretic in nature, the proof relies on (non-abelian and mixed) Hodge theory, and also takes input from the Langlands program, through the work of Esnault-Groechenig \cite{esnault2018cohomologically} 
and Klevdal-Patrikis \cite{klevdalP:g-rigid-local-systems-are-integral}
on integrality of cohomologically rigid local systems.
	
\begin{remark}
\source{canonical-representations-surface-groups}
\notready
	\label{remark:}
\source{canonical-representations-surface-groups}
	Some bound on the rank as in \autoref{theorem:finite-image} is
	necessary. Indeed, there are a number of interesting MCG-finite representations of $\pi_1(\Sigma_{g,n})$ of rank $r\gg g$ with infinite image, as explained in \autoref{question:bounds}. 
	We have no reason to think that the precise bound of
	$\sqrt{g+1}$ is optimal in general. 
	That said the bound cannot be improved too much. It is sharp when $g=1$,
	by \autoref{remark:sharp-g-1}, but not for $g=2,3$, by \autoref{remark:not-sharp-g-2-3}.
	For general $g$ there are MCG-finite representations of $\pi_1(\Sigma_g)$ of rank  
	$2g+1$ with infinite image (\autoref{non-rigid-examples}\eqref{h1-example}), so the bound cannot be improved to be better than linear in $g$.
	We summarize the situation in \autoref{figure:MCG-geography}.
\end{remark}

\begin{figure}
	\centering
	\includegraphics[scale=.7, trim={2cm 18.0cm 0 1cm}]{mcg-geography.pdf}
\caption{The geography of known MCG-finite representations.
\autoref{theorem:finite-image} shows that the region $r<\sqrt{g+1}$ contains
only MCG-finite representations with finite image. As explained in
\autoref{question:bounds}, there are examples of MCG-finite rank $r$ representations with infinite image for $r\gg g$. For $r$ polynomial in $g$, all known examples have virtually solvable image. For $r$ exponential in $g$, there are interesting MCG-finite irreps with infinite image.}\label{figure:MCG-geography}
\source{canonical-representations-surface-groups}
\end{figure}

As explained in ~\cite[Theorem A]{CH:isomonodromic}, 
\autoref{theorem:finite-image} yields a classification result for algebraic solutions to isomonodromy differential equations:
\begin{corollary}
\source{canonical-representations-surface-groups}
\notready
	Let $C$ be a smooth projective curve of genus $g$, and let $D\subset C$
	be a reduced effective divisor. Let $({E},\nabla)$ be a semisimple flat
	vector bundle on $C$ with regular singularities along $D$, with $\rk
	{E}<\sqrt{g+1}$. Suppose that the eigenvalues of the residue matrices of
	$({E},\nabla)$ have real parts in $[0,1)$. Then $({E},\nabla)$ has an algebraic universal isomonodromic deformation if and only if $({E},\nabla)$ has finite monodromy.
\end{corollary}
Note that the condition on residue matrices can always be achieved after a birational gauge transformation, by replacing $({E}, \nabla)$ with the Deligne canonical extension of $({E}, \nabla)|_{C\setminus D}$.

This algebro-geometric reformulation of our main result is not just window
dressing. It is the perspective we will take in the proof of \autoref{theorem:finite-image}. 

\subsection{An application to low rank local systems on $\mathscr M_{g,n}$}

We next record a consequence of our main result for local systems of low rank on families of curves.
As in \autoref{notation:versal-family}, we say a family $\pi:\mathscr C \to \mathscr M$ of smooth proper genus
$g$ curves with geometrically connected fibers, equipped with $n$ disjoint sections $s_i: \mathscr M \to
\mathscr C$, is {\em versal} if the corresponding map $\mathscr M \to \mathscr M_{g,n}$
is dominant. We call 
$\pi^\circ: \mathscr C^\circ := \mathscr C \setminus \bigcup_i s_i(\mathscr M)\to \mathscr
M$ a {\em punctured versal family} of genus $g$ curves.

The following is immediate from \autoref{theorem:finite-image}, by \autoref{proposition:MCG-finite-family-of-curves} below.
\begin{corollary}
\source{canonical-representations-surface-groups}
\notready
	\label{corollary:}
\source{canonical-representations-surface-groups}
	Suppose $\pi^\circ: \mathscr C^\circ \to \mathscr M$
	is a punctured versal family of genus $g$ curves,
	and $\mathbb V$ is a local system on $\mathscr C^\circ$ of rank $<
	\sqrt{g+1}$.
	For any fiber $C^\circ$ of $\pi^\circ$, the local system $\mathbb
	V|_{C^\circ}$ has finite monodromy.
\end{corollary}
%\begin{proof}
%	The restriction $\mathbb V|_{C^\circ}$ corresponds to an MCG-finite
%	representation, since it is fixed by the action of $\pi_1(\mathscr M)$,
%	which has finite index in $\pi_1(\mathscr M_{g,n})$ by the versality
%	assumption, see \autoref{proposition:MCG-finite-family-of-curves}.
%	Then, $\mathbb V|_{C^\circ}$ has finite monodromy by
%	\autoref{theorem:finite-image}.
%\end{proof}

\subsection{The Putman-Wieland conjecture}
\label{subsection:intro-putman-wieland}
\source{canonical-representations-surface-groups}

Among the main ingredients in the proof of \autoref{theorem:finite-image} are some new results toward the Putman-Wieland conjecture, sketched below and
described in detail in \autoref{section:Putman-Wieland}.
%These are also used in the proof of the non-semisimple case of our main result,
%\autoref{theorem:finite-image}.

We now set some notation to state the Putman-Wieland conjecture.
Let $\Sigma_{g,n}$ be a surface and let $\phi: \pi_1(\Sigma_{g,n},
x)\twoheadrightarrow H$ be a surjection onto a finite group, corresponding to
some finite $H$-cover $\Sigma_{g'}\to \Sigma_{g}$ branched over $n$ points. Let $\Gamma$ be the  stabilizer of $\phi$ up to conjugacy in the pure mapping class group of $\Sigma_{g,n+1}$. Then $\Gamma$ acts naturally acts on $H_1(\Sigma_{g'})$. 

We say that {\em the Putman-Wieland conjecture for $(g,n,H)$ holds} if
for every such $H$-cover $\Sigma_{g'} \to \Sigma_g$, 
all nonzero vectors in $H_1(\Sigma_{g'})$ have infinite orbit under $\Gamma$.
The original Putman-Wieland conjecture 
\cite[Conjecture 1.2]{putmanW:abelian-quotients}
asserts that for any fixed $g \geq 2, n \geq 0$,
the Putman-Wieland conjecture for $(g,n,H)$
holds for every $H$.
We prove 
Putman-Wieland conjecture for $(g,n,H)$ for any fixed $H$ once $g$
is sufficiently large. 
In this way, our results are an ``asymptotic" version of the Putman-Wieland conjecture. 
\begin{theorem}
\source{canonical-representations-surface-groups}
\notready\label{theorem:intro-PW}
\source{canonical-representations-surface-groups}
	For $g \geq 2, n \geq 0$, and any finite group $H$ with
	$\#H<g^2$,
	the Putman-Wieland conjecture for $(g,n,H)$ holds.
\end{theorem}
\autoref{theorem:intro-PW} is proven below as a case of \autoref{corollary:asymptotic-putman-wieland}.
More generally, for arbitrary $H$ we prove that the Putman-Wieland conjecture holds for the subspace of $H_1(\Sigma_{g'})$ spanned by irreducible representations of $H$ of small rank.
\begin{theorem}
\source{canonical-representations-surface-groups}
\notready\label{theorem:rho-isotypic-PW-intro}
\source{canonical-representations-surface-groups}
	Let $\rho$ be an irreducible representation of $H$ with $\dim\rho<g$, and let $H_1(\Sigma_{g'})^{\rho}$ be the $\rho$-isotypic component of $H_1(\Sigma_{g'})$. Then no nonzero vector in $H_1(\Sigma_{g'})^{\rho}$ has finite orbit under $\Gamma$.
\end{theorem}
\autoref{theorem:rho-isotypic-PW-intro} is proven below as a special case of
\autoref{theorem:asymptotic-putman-wieland}.

As we shall see in the course of the proof of \autoref{theorem:finite-image},
especially \autoref{lemma:semisimplicity},
the Putman-Wieland conjecture is connected to the classification of MCG-finite representations with virtually solvable image.

\subsection{Arithmetic applications}
\label{subsection:arithmetic}
\source{canonical-representations-surface-groups}
In \autoref{section:arithmetic-applications}, we give a number of applications
to questions in arithmetic geometry. For example, we show in
\autoref{theorem:arithmetic-consequence} that low rank representations of the
arithmetic \'etale fundamental group of a generic curve of genus $g$ have finite
image when restricted to the geometric fundamental group. This verifies a
prediction of the Fontaine-Mazur conjecture, see
\autoref{remark:fontaine-mazur}. 
These results also answer a question of Esnault-Kerz, see \autoref{remark:esnault-kerz}.

As a consequence, we construct many residual
representations of the geometric fundamental group of a generic curve of genus
$g$ which have no lifts to representations of geometric origin. These examples are related to a conjecture of de Jong \cite{de2001conjecture} (proven by Gaitsgory \cite{gaitsgory2007jong}) and a question of Flach, see \autoref{remark:complete-intersection} and \autoref{remark:flach}. 


\subsection{A consequence for free groups}
\label{subsection:free-groups}
\source{canonical-representations-surface-groups}

From \autoref{theorem:finite-image}, we deduce an analogous result for free groups.
\begin{corollary}
\source{canonical-representations-surface-groups}
\notready\label{corollary:free-groups}
\source{canonical-representations-surface-groups}
	Let $F_N$ be a free group on $N$ generators, with $N=2g$ or $N=2g+1$.
	Let $$\rho: F_N\to \on{GL}_r(\mathbb{C})$$ be a representation whose conjugacy class has finite orbit under $\on{Out}(F_N)$. If $r<\sqrt{g+1}$, then $\rho$ has finite image.
\end{corollary}
\begin{proof}
	If $N$ is even, choose an isomorphism $F_N\simeq \pi_1(\Sigma_{g,1})$
	and if $N$ is odd choose $F_N\simeq \pi_1(\Sigma_{g,2})$. Any
	representation of $F_N$ with finite orbit under $\on{Out}(F_N)$ yields a
	representation of a genus $g$ surface group with finite orbit under the
	mapping class group. Such a representation has finite image by \autoref{theorem:finite-image}.
\end{proof}
\begin{remark}
\source{canonical-representations-surface-groups}
\notready
It is natural to ask if the bound on $r$ in \autoref{corollary:free-groups} is sharp. 
There exist (non-semisimple) representations of $F_N$ of fairly low rank ($r =N + 1$) with infinite image and finite orbit under $\on{Out}(F_N)$, see 
\autoref{example:free-group-bounds}, so the bound cannot be improved too much. 
However, in contrast with the case of MCG-finite representations, we do not know any examples of \emph{semisimple} representations of $F_N$ with infinite image and finite orbit under $\on{Out}(F_N)$ as soon as $N\geq 3$.
Indeed, a conjecture of Grunewald and Lubotzky \cite[Conjecture in
\S9.2]{grunewald2009linear}, combined with the main result of
\cite{farb2017moving}, implies that no such examples exist when $N \geq 3$. In the interests of provocation, we conjecture 
(\autoref{conjecture:infinite-image-semisimple})
that such examples \emph{do} exist. 
As evidence, we offer several examples of semisimple MCG-finite representations of surface groups with infinite image, see \autoref{question:bounds}. Note that there are interesting semisimple representations of $F_2$ with finite orbit under $\on{Out}(F_2)$, see \autoref{example:free-group-bounds}.
\end{remark}

\begin{remark}
\source{canonical-representations-surface-groups}
\notready
	\label{remark:explicit-free}
\source{canonical-representations-surface-groups}
One appeal of \autoref{corollary:free-groups} lies in the fact that it admits a completely elementary reformulation, using Nielsen's description of $\on{Aut}(F_N)$,
see \cite{nielsen:die-isomorphismengruppe} 
and also \cite[Theorem 3.2, p. 131]{magnusKS:combinatorial-group-theory}.
Let $(A_1, \cdots, A_N)$ be an $N$-tuple of invertible $r\times r$ complex matrices, with $N=2g$ or $N=2g+1$. Consider the following operations on $N$-tuples of invertible matrices:
\begin{enumerate}
	\item The cyclic permutation $$c: (A_1, A_2, \cdots, A_N)\mapsto (A_2, A_3, \cdots, A_N, A_1)$$
	\item The transposition $$\tau: (A_1, A_2, A_3, \cdots, A_N)\mapsto (A_2, A_1, A_3, \cdots, A_N)$$
	\item The inversion map  $$\epsilon: (A_1, A_2, \cdots, A_N)\mapsto (A_1^{-1},A_2, \cdots, A_N)$$
	\item The Dehn twist $$d: (A_1, A_2, \cdots, A_N)\mapsto (A_1A_2, A_2, \cdots, A_N)$$
\end{enumerate}
	We say $(A_1, \cdots, A_N)$ and $(A_1', \cdots, A_N')$ are \emph{conjugate} if there exists some $B$ such that $A_i=BA_i'B^{-1}$ for all $i$. 
	
	Now suppose that the set of $N$-tuples obtained from $(A_1, \cdots, A_N)$ by repeatedly applying $c, \tau, \epsilon$, and $d$ only intersect finitely many conjugacy classes of $N$-tuples
	$(A_1, \cdots, A_N)$. 
	If $r<\sqrt{g+1}$,
	\autoref{corollary:free-groups} implies that $A_1, \ldots, A_N$ generate a finite subgroup of $\on{GL}_r(\mathbb{C})$. 
\end{remark}
\begin{remark}
\source{canonical-representations-surface-groups}
\notready
	\label{remark:}
\source{canonical-representations-surface-groups}
	Similarly to \autoref{remark:explicit-free}, one could make
	\autoref{theorem:finite-image} explicit, using any one of the known
	generating sets of $\on{Mod}_{g,n}$, but the formulas are a bit more
	involved, as in \cite[\S6]{CH:isomonodromic}. It would in our view be of great interest to find a \emph{proof} of \autoref{theorem:finite-image} or \autoref{corollary:free-groups} of a similarly explicit nature.
\end{remark}
	
\begin{remark}
\source{canonical-representations-surface-groups}
\notready
A result analogous to \autoref{corollary:free-groups} follows immediately for
characteristic quotients of free or surface groups. For example, let $G$ be a
characteristic quotient of a free group $F_N$ on $N=2g$ or $N=2g+1$ generators,
and $\rho: G\to \on{GL}_r(\mathbb{C})$ a representation whose conjugacy class has finite orbit under $\on{Out}(G)$. If $r<\sqrt{g+1}$, then $\rho$ has finite image by \autoref{corollary:free-groups}.
\end{remark}


\subsection{Cohomological results}
\label{subsection:cohomological}
\source{canonical-representations-surface-groups}
One of the new technical inputs we use to achieve \autoref{theorem:finite-image} is
an analysis of the cohomology of unitary local systems on families of curves.
Our main result, which follows from an analysis of the derivative of the period
map associated to the mixed Hodge structure on the cohomology of unitary local
systems, is the following.
We adapt notation described later in \autoref{notation:versal-family}.
\begin{theorem}
\source{canonical-representations-surface-groups}
\notready\label{theorem:rank-bound-subsystem}
\source{canonical-representations-surface-groups}
Let $\pi:\mathscr{C}\to \mathscr{M}$ be a smooth proper family of $n$-pointed
curves of genus $g$ with geometrically connected fibers, so that the associated
map $\mathscr{M}\to \mathscr{M}_{g,n}$ is dominant \'etale. Let $\pi^\circ:
\mathscr{C}^\circ\to \mathscr{M}$ be the associated family of punctured curves.
Let $\mathbb{V}$ be a unitary local system on $\mathscr{C}^\circ$. Then any
nonzero sub-local system of $R^1\pi_*^\circ \mathbb{V}$ has rank at least $2g-2\rk\mathbb{V}$. 	
\end{theorem}
We prove this in \autoref{subsection:rank-bound-proof}, and give a related cohomological vanishing theorem with milder unitarity hypotheses in \autoref{theorem:unitary-rigid}.
Note that \autoref{theorem:rank-bound-subsystem} shows that $H^0(\mathscr{M}, R^1\pi_*^\circ \mathbb{V})=0$ if
$g>\on{rk}\mathbb{V}$. 
This result is used several times in the proof of
\autoref{theorem:finite-image}, as described in
\autoref{subsection:outline-of-proof}.
It is especially interesting when $\mathbb{V}$ has finite monodromy, where it is closely related to the Putman-Wieland conjecture, as described
above in \autoref{subsection:intro-putman-wieland} and in more detail in \autoref{section:Putman-Wieland}.


\subsection{Prior work}
\label{subsection:prior-work}
\source{canonical-representations-surface-groups}
\subsubsection{MCG-finite representations and Painlev\'e VI}
\label{subsubsection:prior-mcg-finite}
\source{canonical-representations-surface-groups}
As mentioned in \autoref{subsection:overview}, there has been a huge amount of effort put forth towards classifying algebraic solutions to the Painlev\'e VI equation, arguably beginning with work of Riemann \cite{riemann1857beitrage} and Schwarz \cite{schwarz1873ueber}
and culminating in the classification \cite{lisovyy2014algebraic}. 
Calligaris-Mazzocco \cite{calligaris2018finite} classified algebraic
solutions satisfying several conditions in the case $g=0, n=5, r=2$, but the complete classification in
this case remains open. See also \cite{michel-hurwitz} for a related result in the genus zero case, when the local monodromy matrices are given by reflections; this appears to be the first result showing that MCG-finite representations (satisfying certain conditions) have finite image.

For all the work invested in the algebraic solutions to Painlev\'e VI, this only
addresses the case of classifying  MCG-finite representations with $g = 0, n = 4,
r = 2$, and trivial determinant, as explained in \cite{lisovyy2014algebraic}.
In higher genus, all work that we know of has been on the case of
$2$-dimensional representations.
The beautiful paper \cite{biswas2017surface} handles the $2$-dimensional case ($r = 2, g \geq 1, n \geq 0$) with trivial
determinant.
The non-semisimple case $r = 2, g \geq 1, n \geq 0$ with arbitrary determinant
is explained in \cite[Theorem B]{CH:isomonodromic}.
These results appear to depend crucially on the assumption $r=2$.
%Note that any representation with finite image is MCG-finite, and in fact
%when $g \geq 2$ and $r = 2$, all such representations were shown to have finite
%image.

\begin{remark}
\source{canonical-representations-surface-groups}
\notready
	\label{remark:}
\source{canonical-representations-surface-groups}
	The above-mentioned results in rank $2$ led Junho Peter Whang to ask whether, for $g\gg r$, all MCG-finite rank $r$ representations of $\Sigma_{g,n}$
have finite image \cite[Question 1.5.3]{lawrence2019representations}.
More generally, motivated by the $p$-curvature conjecture, Mark Kisin asked whether MCG-finite representations
necessarily have finite image,
see \cite[p. 3]{biswas2017surface} and \cite[p. 1]{sinz:thesis}. (Note that
counterexamples are known in general, see ~\cite[Theorem
5.1]{biswas2018representations} for counterexamples coming from TQFT techniques or 
\cite[Example 3.3.1]{lawrence2019representations} for counterexamples coming
from the Kodaira-Parshin trick, as well as
\autoref{question:bounds} of this paper.)
In this way, \autoref{theorem:finite-image} answers Whang's question
affirmatively, and provides a positive answer to Kisin's question in the regime
$g > r^2-1$.
\end{remark}

There has also been much work put towards finding interesting \emph{examples} of MCG-finite representations, notably \cite{doran:isomonodromic, diarra2013construction, girand:a-new-two-parameter-family, diarra2015ramified, girand2016algebraic}. 

In a more arithmetic direction, \cite{bourgainGS:announcement} 
announced striking results on strong approximation for Markoff triples, obtained by analyzing the
arithmetic properties of subvarieties of the character variety parametrizing $2$-dimensional
representations of $\pi_1(\Sigma_{1,1})$.
The first step in their analysis is to determine the finite orbits of the
mapping class group action on this character variety \cite[last paragraph of p.
132]{bourgainGS:announcement}.
In this way, our \autoref{theorem:finite-image} can be viewed as a necessary
first step 
toward attempting to generalize their approach to 
higher rank, higher genus character
varieties.

\subsubsection{MCG-finite representations and geometric topology}

In the world of low-dimensional topology a number of authors have studied the
dynamics of mapping class group actions on character varieties; finite orbits
are the same as semisimple MCG-finite representations. Kasahara
\cite{kasahara2015visualization} relates fixed points of this action
corresponding to faithful representations to the well-known question of
linearity of the mapping class group. Goldman and many other authors have
studied ergodicity of these actions, see \cite{goldmanmapping} and the
references therein. Previte-Xia \cite{previte2000topological,
previte2002topological} study the relationship between density of the image of
an $SU(2)$-representation and density of its mapping class group orbit. In some
sense our main result is a partial answer to \cite[Question
2.7]{goldmanmapping}, which asks for necessary and sufficient conditions for a
representation to have dense orbit under the mapping class group. We
characterize the most extreme possible failure of density, namely the case of
finite orbit.

There are also a number of related results on representations of the mapping class group.
Farb, Lubotzky, and Minsky \cite[Theorem 1.6]{farbLM:rank-1-phenomena} show there are no faithful linear
representations of finite index subgroups of $\on{Mod}_{g,0}$ of dimension $< 2 \sqrt{g-1}$. See \autoref{remark:FLM-rank-1-result} for a comparison of our results to theirs.
See also
\cite{franksH:triviality,korkmaz:low-dimensional-linear-representations,funar:two-questions, kielak-pierro}
for bounds on the dimension of representations of 
$\on{Mod}_{g,n}$, although these results only address representations of the full mapping
class group, as opposed to representations of finite index subgroups.
%There has also been significant work bounding the dimension of representations
%of $\on{Mod}_{g,n}$ with infinite image.
%Of course, there is a representation of dimension $2g$ where 
%$\on{Mod}_{g,n}$ acts on $H_1(\Sigma_{g,n}, \mathbb C)$.
%However, in 
%\cite{franksH:triviality}, it is shown there any representation of 
%$\on{Mod}_{g,n}$ of dimension $<2g$ is trivial.
%This was also proven in \cite{korkmaz:low-dimensional-linear-representations}
%and improves the weaker bound from \cite{funar:two-questions}.
%However, the above results only address representations of the full mapping
%class group, whereas we are investigating properties of finite index subgroups
%of the mapping class group, which tend to be much more subtle.



\subsubsection{The Putman-Wieland conjecture} The other main contribution of
\label{subsubsection:prior-putman-wieland}
\source{canonical-representations-surface-groups}
this paper, towards the Putman-Wieland conjecture, has a number of precursors,
notably \cite{looijenga:prym-representations,grunewald2015arithmetic, looijenga2021arithmetic}.
Our Hodge-theoretic approach is related to an approach suggested by Looijenga \cite{looiejenga:AG}. 
The strongest previous result towards the Putman-Wieland conjecture is perhaps 
the main result of \cite{grunewald2015arithmetic}, which says that for certain quotients of $\pi_1(\Sigma_g)$, 
called \emph{$\phi$-redundant} quotients,
much more than the Putman-Wieland conjecture is true --- the monodromy representations considered by the conjecture have very large image, commensurable with an arithmetic group. There has also been interesting recent work related to the Putman-Wieland conjecture by Markovi\'c and Markovi\'c-To\u{s}i\'c \cite{markovic, markovic2}.

The analogue of the Putman-Wieland conjecture for graphs, as opposed to surfaces, has been proven by Farb and Hensel \cite{farb2017moving}.


\subsection{Outline of the proof}\label{subsection:outline-of-proof}
\source{canonical-representations-surface-groups}
We now sketch the proof of \autoref{theorem:finite-image}, which is loosely inspired by Katz's proof of the $p$-curvature conjecture for the Gauss-Manin connection \cite{katz:pcurvature-and-hodge}.
Following this, we sketch the proof of the vanishing results we use (\autoref{theorem:rank-bound-subsystem} and its consequence \autoref{theorem:unitary-rigid}) in
\autoref{subsubsection:rigidity-idea}, which are used in different ways in each of the three steps listed below of the proof of \autoref{theorem:finite-image}.
A schematic diagram outlining the main elements of the proof is depicted in
\autoref{figure:proof-schematic}.

\begin{figure}
	\centering
\adjustbox{scale=.75,center}{
   \begin{tikzcd}[column sep = 2.3em]
	\qquad &&
\text{Lem.}~\ref{lemma:gl-defined-over-number-field} \ar{d}&
  \text{Lem.}~\ref{lemma:defined-over-number-field} \ar{l} \\
%5
  \qquad &&
  \underbrace{\text{Prop.}~\ref{proposition:unitary-implies-finite}}_{\text{Unitary
  case}}
  \ar{dl}{\text{\cite[Non-abelian Hodge theory]{mochizuki2006kobayashi}}} 
  \ar[swap]{dl}{\text{\cite[Isomonodromy]{LL:geometric-local-systems}}}  &
  \text{Prop.}~\ref{proposition:cohomologically-rigid} 
  \ar{l}{\text{\cite{mochizuki2006kobayashi}}} 
  \ar[swap]{l}{\text{\cite{LL:geometric-local-systems}}}
  \ar[swap]{u}{\text{\cite[Langlands]{esnault2018cohomologically, klevdalP:g-rigid-local-systems-are-integral}}} &
  \text{Prop.}~\ref{proposition:generic-parabolic-global-generation} \ar{d} \\
%6
  \text{Thm.}~\ref{theorem:finite-image} &
  \underbrace{\text{Thm.}~\ref{theorem:finite-image-semisimple}}_{\text{Semisimple
  case}}\ar{l}&
  \text{Lem.}~\ref{lemma:global-rigidity} \ar{l} &
  \underbrace{\text{Thm.}~\ref{theorem:unitary-rigid}}_{\text{Cohomological
  vanishing}} \ar{l} \ar{u}&
  \text{Prop.}~\ref{proposition:pairing-result} \ar{dl}{\text{Mixed Hodge
  theory and \cite{mehta1980moduli}}} \\
%7
  \qquad & \text{Lem.}~\ref{lemma:semisimplicity} \ar{ul}&
  \underbrace{\text{Thm.}~\ref{theorem:asymptotic-putman-wieland}}_{\text{Putman-Wieland}} \ar{l}&
	  \underbrace{\text{Thm.}~\ref{theorem:rank-bound-subsystem}}_{\text{Cohomological
	  rank bound}}\ar{l} \ar{u} 
	  &
	  \underbrace{\text{Thm.}~\ref{theorem:period-map-computation}}_{\text{Period map}}
	  \ar{l}
	  \end{tikzcd}
}
\caption{
A diagram depicting the structure of the proof of
the main result,
\autoref{theorem:finite-image}, in the shape of a boat.}
\label{figure:proof-schematic}
\source{canonical-representations-surface-groups}
\end{figure}


\subsubsection{Step 1. The unitary case.} 
\label{subsubsection:step-1}
\source{canonical-representations-surface-groups}
Every unitary representation is a direct sum of irreducible unitary representations.
Therefore, we suppose $\rho$ is unitary, irreducible, and MCG-finite, 
with $\rk\rho<\sqrt{g+1}$. We construct from $\rho$ a finite
\'etale cover $\mathscr{M}$ of $\mathscr{M}_{g,n}$, with associated family of
punctured curves $\pi^\circ:\mathscr{C}^\circ\to \mathscr{M}$, and a projective unitary local
system $\mathbb{V}$ on $\mathscr{C}^\circ$ whose restriction to a fiber
$C^\circ$ of
$\pi^\circ$ has monodromy given by $\rho$. 
Applying \autoref{proposition:cohomologically-rigid}, (which is a fairly 
straightforward consequence of our cohomological vanishing result,
\autoref{theorem:unitary-rigid}, applied to $\on{ad}\mathbb{V}$), 
shows that $\mathbb{V}$ is cohomologically rigid. 
%Note that this only requires the restriction of $\mathbb V$ to a fiber of $\pi$
%to be unitary.
The main result of \cite{klevdalP:g-rigid-local-systems-are-integral} then gives that $\rho$ is defined over the ring of integers $\mathscr{O}_K$ of some number field $K$.
Moreover, using \autoref{proposition:gl-lift},
by replacing $\mathscr M$ by a dominant \'etale scheme over it, along which
certain cohomological lifting obstructions vanish, we can assume
$\mathbb V$ lifts from a projective local system to a bona fide local system.
%riga

By compactness of the unitary group and discreteness of $\mathscr{O}_K$, it
suffices to show that for each embedding $\iota:
\mathscr{O}_K\hookrightarrow\mathbb{C}$, $\rho\otimes_{\mathscr{O}_K, \iota}
\mathbb{C}$ is unitary. (We know this unitarity for one such $\iota$ by assumption, but not
the others.) The rigidity of $\mathbb{V}$ implies by non-abelian Hodge theory
that these $\rho\otimes_{\mathscr{O}_K, \iota} \mathbb{C}$ underlie complex polarizable variations of Hodge structure for any complex structure on $\Sigma_{g,n}$. 
Hence, given their low rank, these local systems are unitary by \cite[Theorem
1.2.12]{LL:geometric-local-systems}.

\subsubsection{Step 2. The semisimple case.} 
\label{subsubsection:step-2}
\source{canonical-representations-surface-groups}
Now suppose $\rho$ is an arbitrary
semisimple, MCG-finite representation with $\rk\rho<\sqrt{g+1}$. Again, we
associate to $\rho$ a local system $\mathbb{V}$ on a family of curves
$\pi^\circ:\mathscr{C}^\circ\to \mathscr{M}$, so that $\mathscr M$ has a
dominant \'etale map to $\mathscr M_{g,n}$ and whose fibral monodromy is given by $\rho$.
By non-abelian Hodge theory, we may deform $\mathbb{V}$ to a local system
$\mathbb V_0$ underlying a complex polarizable
variation of Hodge structure. By
\cite[Theorem 1.2.12]{LL:geometric-local-systems}, $\mathbb{V}_0$ has unitary monodromy when
restricted to a fiber of $\pi^\circ$. By the unitary case,
\autoref{subsubsection:step-1},
$\mathbb{V}_0$ thus has finite monodromy when restricted to a fiber of $\pi^{\circ}$. Note that even if $n=0$, i.e.~$\pi^\circ$ is proper, we here need to use non-abelian Hodge theory for non-proper varieties, as the total space $\mathscr{C}^\circ$ will not be proper.

It remains to argue that the restriction $\mathbb V_0|_{C^\circ}$ of
$\mathbb V_0$ to a fiber $C^\circ$ of
$\pi^\circ$ agrees with the restriction $\mathbb V|_{C^\circ}$,
corresponding to $\rho$.
Recall that 
$\mathbb V_0$ and $\mathbb V$ were only deformation equivalent, so it may be
surprising that they necessarily restrict to the same local system on fibers.
We verify this agreement 
in \autoref{lemma:global-rigidity}
through another application of
\autoref{theorem:unitary-rigid}, which 
which tells us that since the fibral monodromy $\mathbb{V}_0|_{C^\circ}$ 
is unitary of low rank, it
does not admit non-trivial MCG-finite deformations. 

Note that $\mathbb{V}_0$ may not a priori be unitary; we only know that $\mathbb V_0|_{C^\circ}$ is unitary. In particular, it is not clear whether $\mathbb{V}_0$ is necessarily cohomologically rigid.

\subsubsection{Step 3. The general case} 
\label{subsubsection:step-3}
\source{canonical-representations-surface-groups}
The crucial input for dealing with
non-semisimple representations is our work towards the
Putman-Wieland conjecture. By the above it is enough to show that low-rank
MCG-finite representations are semisimple, i.e.~we wish to verify that certain
extensions split. 
For simplicity, let's suppose for the purpose of this sketch that $\rho$ is an extension of two irreducible MCG-finite representations
$\rho_1$ and $\rho_2$. 
By the previous step, $\rho_1$ and $\rho_2$ have finite monodromy, so after
passing to a finite cover $\Sigma_{g',n'}$ of $\Sigma_{g,n}$, we may assume
$\rho_1$ and $\rho_2$ have trivial monodromy.
The splitting of this extension of $\rho_2$ by $\rho_1$ corresponds to the vanishing of a certain
element in $\on{Ext}^1_{\pi_1(\Sigma_{g',n'})}(\rho_1, \rho_2)$.
Because we arranged that $\rho_1$ and $\rho_2$ have trivial monodromy
on $\pi_1(\Sigma_{g',n'})$, the above extension class corresponds to a map
$\pi_1(\Sigma_{g',n'}) \to \rho_1^\vee \otimes \rho_2$, with unipotent abelian
image. Hence, it factors through $H_1(\Sigma_{g',n'})$,
and so defines a low rank subspace of $H^1(\Sigma_{g',n'})$, stable under a
finite index subgroup of $\on{Mod}_{g,n+1}$.
We verify in \autoref{lemma:semisimplicity}, 
using our main result toward the Putman-Wieland conjecture,
\autoref{theorem:asymptotic-putman-wieland},
that such a subspace cannot exist. Theorem \autoref{theorem:asymptotic-putman-wieland} itself follows more or less immediately from \autoref{theorem:rank-bound-subsystem}.

%It remains to explain the proof of 
%\autoref{theorem:asymptotic-putman-wieland}.
%From the above data, one can construct a family of punctured curves $\mathscr C^\circ$ over $\mathscr M$,
%for $\mathscr M$ a scheme with a dominant \'etale map to $\mathscr M_{g,n}$.
%Moreover, one obtains finite, (hence unitary,) local systems $\mathbb U_\sigma$ on $\mathscr C^\circ$, for $\sigma$ varying over the
%irreducible representations of the Galois group of the cover $\Sigma_{g',n'} \to
%\Sigma_{g,n}$.
%Then, the low rank subspace of $H^1(\Sigma_{g',n'})$ above yields a low rank
%sub-local system of $R^1 \pi^\circ_* \mathbb U_\sigma$, which contradicts 
%\autoref{theorem:rank-bound-subsystem}.
%their failure to do so would yield low-rank sub-local
%systems of $R^1\pi_*^\circ\mathbb{V}$ for $\pi^\circ: \mathscr{C}^\circ\to
%\mathscr{M}$ as above, for some $\mathbb{V}$ with finite monodromy. 
%This contradicts
%\autoref{theorem:rank-bound-subsystem}.

\subsubsection{The proof of our cohomological vanishing results}
\label{subsubsection:rigidity-idea}
\source{canonical-representations-surface-groups}
Note that we used \autoref{theorem:rank-bound-subsystem}, or its immediate consequence,
\autoref{theorem:unitary-rigid},
in every step above.
Before explaining the idea of the proof of these results, let us recall the
setting. We begin with a scheme $\mathscr M$ with a dominant \'etale map
$\mathscr M \to \mathscr M_{g,n}$.  We denote the associated family of punctured curves by $\mathscr
C^\circ \to \mathscr M$. We are given a local system $\mathbb V$ on $\mathscr C^\circ$.
\autoref{theorem:rank-bound-subsystem} says that if $\mathbb{V}$ is unitary,
there are no sub-local systems of $R^1 \pi^\circ_* \mathbb V$ of low rank (that is, rank less
than $2g - 2\rk \mathbb V$).
Let $C^\circ$ be a fiber of $\pi^\circ$.
%Now, let $m \in \mathscr M$ be a point and $C^\circ := \mathscr C^\circ_m$
%denote a fiber of $\pi^\circ$.
We deduce \autoref{theorem:unitary-rigid}, which says that when $\mathbb
V|_{C^\circ}$
is unitary and has low rank (less than $g$), then $R^1 \pi^\circ_* \mathbb V$ has no global sections.
The benefit of \autoref{theorem:unitary-rigid} is that we only need check
unitarity on fibers, but the cost is that we are only able to rule out global
sections, i.e. trivial sub-local systems, instead of arbitrary low rank
sub-local systems.

First, we deduce \autoref{theorem:unitary-rigid} from
\autoref{theorem:rank-bound-subsystem}.
The idea is to use the assumption that
$\mathbb V$ is unitary on fibers to
reduce via \autoref{lemma:unitary-direct-sum} to the case that $\mathbb V$ is a
tensor product $\mathbb U \otimes (\pi^\circ)^* \mathbb W$, with $\mathbb{U}$ unitary. 
A nonzero global section of $R^1\pi^\circ_*(\mathbb U \otimes (\pi^\circ)^* \mathbb W)=(R^1\pi^\circ_*\mathbb{U})\otimes \mathbb{W}$ yields a nonzero map $\mathbb W^\vee \to R^1 \pi^\circ_* \mathbb U$, and hence a
low rank sub-local system of $R^1 \pi^\circ_* \mathbb U$, contradicting 
\autoref{theorem:rank-bound-subsystem}.

The proof of
\autoref{theorem:rank-bound-subsystem} boils down to an analysis of the derivative of the period map associated to the complex variation of mixed Hodge structure on $R^1\pi^\circ_*\mathbb{V}$, for $\mathbb{V}$ unitary.
%The proof is based on an analysis of
%derivatives of period maps associated to the mixed Hodge structure on the
%cohomology of unitary local systems. 
%If we had a low rank sub-local system of $R^1 \pi^\circ_* \mathbb V$,
%we would, in particular, obtain such a sub-local system over a first order
%neighborhood of any point $m \in \mathscr M$.
%More precisely, let $C^\circ$ is the fiber of $\pi^\circ$ over $m$ and $C$ is the corresponding
%proper curve, so that $C^\circ = C - D$, let $E$ denote the Deligne canonical extension to $C$ of the vector bundle
%$\mathbb V|_{C^\circ} \otimes \mathscr O_{C^\circ}$.
%Then, a sub-local system as above would yield a subspace $H^1(C^\circ, \mathbb
%V|_{C^\circ})$ which does not vary to first order over $\mathscr M$. 
%Using mixed Hodge theory to analyze the associated period map
%$\mu: H^0(C, E \otimes \omega_C(D)) \to H^1(C, E) \otimes \Omega_{\mathscr M}$,
%we obtain an element $v \in H^0(C, E \otimes \omega_C(D))$ so that $\mu(v)$ has
%low rank.
We identify this derivative with a natural multiplication map $$H^0(E\otimes \omega_C(D))\otimes H^0(E^\vee\otimes \omega_C)\to H^0(\omega_C^{\otimes 2}(D))$$ in
\autoref{theorem:period-map-computation}, where $C$ is the smooth projective compactification of $C^\circ$ with boundary $D$, and $E$ is the vector bundle on $C$ associated to $\mathbb{V}|_{C^\circ}$ by the Mehta-Seshadri correspondence \cite{mehta1980moduli}. By the theorem of the fixed part, the existence of a low rank sub-local system of $R^1\pi_*^\circ\mathbb{V}$ places restrictions on this map, which we rule out
%We then show no such element $v$ can exist in \autoref{proposition:pairing-result}
%by studying the pairing $(E \otimes \omega_C(D)) \otimes (E^\vee\otimes
%\omega_C) \to \omega_C^{\otimes 2}(D)$ 
%associated to the above multiplication map.
 using vector bundle methods,
as developed in \cite[\S6]{LL:geometric-local-systems},
ultimately relying on Clifford's theorem for vector bundles.

%\subsection{Structure of the paper}
%\label{subsection:structure}
%In \autoref{section:representation-preliminaries}, we describe some geometric
%and representation-theoretic preliminaries for local systems on families of
%curves, and a number of constructions we will use throughout the paper. In
%\autoref{section:Hodge-theoretic-preliminaries}, we describe several results on
%the mixed Hodge structure on the cohomology of a unitary local system, and we
%recall the facts from non-abelian Hodge theory that we will use. In
%\autoref{section:unitary-period-map}, we compute the derivative of the
%associated period map (with proofs in \autoref{appendix}).
%\autoref{section:vanishing} contains the proof of the main cohomological result, \autoref{theorem:rank-bound-subsystem}, and its consequences. We give applications to the Putman-Wieland conjecture in \autoref{section:Putman-Wieland}. The proof of the main result, \autoref{theorem:finite-image}, is given in \autoref{section:main-proof}. In \autoref{section:arithmetic-applications}, we give a number of arithmetic applications. Finally, \autoref{section:questions} contains several examples, questions, and conjectures motivated by our results.

\subsection{Notation}
\label{subsection:notation}
\source{canonical-representations-surface-groups}
Unless otherwise stated, we will work over the field of complex numbers $\mathbb{C}$. 
In particular, we will lift this restriction in
\autoref{section:arithmetic-applications}.

\begin{notation}
\source{canonical-representations-surface-groups}
\notready
	\label{notation:versal-family}
\source{canonical-representations-surface-groups}
	We fix non-negative integers $(g,n)$ so that $n \geq 1$ if $g = 1$ and
	$n \geq 3$ if $g = 0$, i.e., $\Sigma_{g,n}$ is hyperbolic.
	Let $\mathscr{M}$ be a connected complex variety. A \emph{family
	of $n$-pointed curves of genus $g$ over $\mathscr{M}$} is a smooth
	proper morphism $\pi: \mathscr{C}\to \mathscr{M}$ of relative dimension one, with geometrically
	connected genus $g$ fibers, equipped with $n$ sections $s_1,\cdots,
	s_n: \mathscr{M}\to \mathscr{C}$ with disjoint images. Call such a family is \emph{versal} if the
	induced map $\mathscr{M}\to \mathscr{M}_{g,n}$ is dominant and \'etale.
	Here $\mathscr M_{g,n}$ denotes the Deligne-Mumford moduli stack of $n$-pointed genus
	$g$ smooth proper curves with geometrically connected fibers.

	If $\pi: \mathscr{C}\to \mathscr{M}$ is a family of $n$-pointed curves, we let
	$\mathscr D := \coprod_{i=1}^n \on{im}(s_i)$ denote the images of
	the sections, which is finite \'etale of degree $n$ over $\mathscr M$.
	Also let
	$\mathscr{C}^\circ:=\mathscr{C}\setminus \bigcup_i \on{im}(s_i)$,
	let $j: \mathscr C^\circ \hookrightarrow \mathscr{C}$ be the natural
	inclusion, and let $\pi^\circ := \pi \circ j : \mathscr C^\circ \to \mathscr	M$ denote the composition. We will refer to $\pi^\circ: \mathscr{C}^\circ\to \mathscr{M}$ as the \emph{associated family of punctured curves}. If $\pi^\circ$ arises as the family of punctured curves associated to a versal family of $n$-pointed curves, we will call it a \emph{punctured versal family}.
	We will frequently use $m \in \mathscr M$ as a basepoint, and $c \in
	\mathscr C^\circ$ as a basepoint with $\pi^\circ(c) = m$.
	We use $C^\circ$ as notation to denote the fiber $(\pi^\circ)^{-1}(m) = \mathscr
	C^\circ_m$.
	\begin{equation}
		\label{equation:}
\source{canonical-representations-surface-groups}
		\begin{tikzcd}
			\mathscr C^\circ \ar {r}{j} \ar {rd}{\pi^\circ} &
			\mathscr C \ar {d}{\pi} & \mathscr D \ar{l} \\
			& \mathscr M \ar{ur}{s_1} \ar[bend
			right=70,swap,ur,"s_n"]
			\arrow[bend right = 60, ur, draw=none, "\ddots"]
		 &
		\end{tikzcd}
	\end{equation}
\end{notation}

%\begin{notation}[Conventions on basepoints]
%	\label{notation:basepoint}
%	We will use the following conventions for basepoints.
%We use $m$ for a basepoint of $\mathscr M$ and $c$ for a basepoint of $\mathscr
%C^\circ$ lying over $m$.
%Throughout we shall introduce auxiliary schemes
%$\mathscr M', \mathscr M_i,\mathscr C'^\circ, \mathscr C_i^\circ$ 
%with maps $\mathscr M' \to \mathscr M, \mathscr M_i \to \mathscr M, \mathscr
%C'^\circ
%\to \mathscr C, \mathscr C_i^\circ \to \mathscr C$
%and similarly use $m', m_i, c', c_i$ for basepoints of these schemes.
%We shall assume (adjusting $m$ and $c$ if necessary) that $m' \mapsto m, m_i
%\mapsto m, c' \mapsto c, c_i\mapsto c$.
%\end{notation}

\begin{notation}
\source{canonical-representations-surface-groups}
\notready
	\label{notation:adjoint}
\source{canonical-representations-surface-groups}
	For $G$ an algebraic group with derived subgroup $G^{\on{der}}$ and
corresponding Lie algebra $\mathfrak g^{\on{der}}$, we use
$\on{Ad} : G \to \on{GL}(\mathfrak g^{\on{der}})$
to denote the natural action of $G$ on $\mathfrak{g}^{\on{der}}$ by conjugation.
Given a representation $\rho: \pi_1(X, x) \to G(\mathbb C)$, let $\on{ad}(\rho)
:= \rho \circ \on{Ad}: \pi_1(X, x) \to \on{GL}(\mathfrak g^{\on{der}})$. 
In particular, given
$\rho : \pi_1(X, x) \to \on{GL}_r(\mathbb C)$ or 
$\rho : \pi_1(X, x) \to \on{PGL}_r(\mathbb C)$
we use $\on{ad}(\rho)$ to denote the composite map
$\on{ad}(\rho) : \pi_1(X, x) \to \on{GL}(\mathfrak {pgl}_r(\mathbb C))$.
Under the identification between local systems on a connected space and representations of the
fundamental group, if $\mathbb W$ is a local system on $X$ corresponding to some representation $\rho$, 
we use $\on{ad}\mathbb W$ to denote the local system corresponding to
$\on{ad}\rho$.

A representation $\pi_1(X,x)\to \on{GL}_r(\mathbb{C})$ is \emph{unitary} if its image has
compact closure. 
A complex local system is unitary if its monodromy representation is unitary. By
an averaging argument, unitary representations are exactly those that preserve a positive-definite Hermitian form on $\mathbb{C}^n$, i.e.~they are conjugate to a representation that factors through the unitary group $\on{U}(n)$.
\end{notation}

\begin{notation}
\source{canonical-representations-surface-groups}
\notready
For a pointed finite-type scheme or Deligne-Mumford stack $(X, x)$ over
$\mathbb{C}$, we will use $\pi_1(X, x)$ to denote the topological fundamental group 
 of the associated complex-analytic space or analytic stack.
Similarly, for a local system $\Lambda$ on $X$, we use
$H^i(X, \Lambda)$ to denote the singular cohomology of the associated complex-analytic space or stack, unless otherwise stated. We will use $\on{Mod}_{g,n}$ to denote the mapping class group of an orientable surface of genus $g$ with $n$ punctures, and $\on{PMod}_{g,n}=\pi_1(\mathscr{M}_{g,n})$ 
(see \autoref{lemma:mgn-fundamental-group})
to denote pure mapping class group, i.e.~the subgroup of $\on{Mod}_{g,n}$ of index $n!$ preserving the punctures pointwise.
\end{notation}

\subsection{Acknowledgements}

We would like to thank H\'el\`ene Esnault
for extremely helpful comments, and especially for finding and helping us correct a substantial gap in the original proof of \autoref{theorem:finite-image-semisimple}. We thank Andrew Putman for encouraging us in this project, and Ian Agol and Dawid Kielak for bringing the Putman-Wieland conjecture to our attention.
We also thank an anonymous referee,
Philip Boalch,
Anand Deopurkar,
Charles Doran,
Benson Farb,
Matt Kerr,
Mark Kisin,
Christian Klevdal,
Joshua Lam,
Brian Lawrence,
Alexander Lubotzky,
Vladimir Markovi\'c,
Takuro Mochizuki,
Stefan Patrikis,
Alexander Petrov,
Mihnea Popa,
Will Sawin,
Carlos Simpson,
Andrew Snowden,
Junho Peter Whang,
Ben Wieland,
Melanie Matchett Wood, 
and Kang Zuo.
Landesman was supported by the National Science Foundation under Award No. DMS-2102955 
and Litt was supported by NSF grant DMS-2001196.
This material is based upon work supported by the Swedish Research Council under grant no. 2016-06596 while the authors were in residence at Institut Mittag-Leffler in Djursholm, Sweden during the fall of 2021.
